\newpage
\section{MVP}

\subsection{Requisitos Priorizados}
\par A partir das discussões na reunião pós entrega do Sprint 2, pudemos ter uma visão melhor
do objetivo dos usuários, e assim, conseguimos decidir quais requisitos e funções priozirar 
para a implementação de nosso MVP
\par Entendemos que o principal objetivo da plataforma a ser desenvolvida é promover a interação
social entre os seniors. Portanto o foco estará nas reuniões ao vivo e tudo que é necessário para
que elas possam ser aplicadas, como calendário e inscrissões.
\par Segue a lista de requisitos priorizados na implementação, você pode voltar para a seção 
\ref{sec:requisitos} para visualizar todos os requisitos e histórias de usuário na integra.

\begin{enumerate}
    \item RF1 - Visualizar atividades disponíveis
    \item RF2 - Inscrever-se em atividade
    \item RF3 - Atividades online
    \item RF6 - Lembretes de atividades
    \item RF9 - Criação e edição de atividades
    \item RF11 - Gerenciamento de lista de participantes
\end{enumerate}

\par Como o mvp já é desenvolvido para o uso pelo público 65+, foi dada a prioridade aos requisitos
não funcionais de usabilidade, acessibilidade e segurança. Desempenho e compatibilidade serão trabalhados
mas com menos prioridade.
